% Môn: Vật lý
% Lớp: 12
% Chương: Chương II. Khí lí tưởng
% Bài: L12C2 Bài 8. Mô hình động học phân tử chất khí · Mô hình động học phân tử về cấu tạo chất và chuyển động Brown
% ID: L12C1B1-TN-42
% Mức: NB
% ID: L12C1B1-91-TN
% ID: L12C1B1-91-TN
% Mức: NB
% ID: L12C2B8-61-TN
%=== Câu 1 ===%
\dangbt{Mô hình động học phân tử về cấu tạo chất và chuyển động Brown}
\begin{ex}
% ID: L12C2B8-61-TN
% Mức: NB
	{\color{red}\footnotesize\ttfamily [L12C1B1-TN-42]}\par %HienID
	Chuyển động Brown là
	\choice
	{chuyển động có quỹ đạo xác định của các phân tử chất lỏng}
	{\True chuyển động hỗn loạn, không ngừng của các hạt nhỏ (như hạt phấn hoa) lơ lửng trong chất lỏng hoặc chất khí, do bị các phân tử môi trường va chạm}
	{sự rung động tại chỗ của nguyên tử trong mạng tinh thể chất rắn}
	{sự khuếch tán có quy luật của các phân tử khí trong chất lỏng}
	\loigiai{
		Chuyển động Brown là chuyển động hỗn loạn, không ngừng của các hạt rất nhỏ lơ lửng trong chất lỏng hoặc chất khí, gây ra bởi sự va chạm không cân bằng liên tục của các phân tử môi trường xung quanh vào hạt đó.
		Kết luận: Phương án đúng là ý (2).
	}
\end{ex}

% ID: L12C1B1-TN-43
% Mức: TH
% ID: L12C1B1-92-TN
% ID: L12C1B1-92-TN
% Mức: TH
% ID: L12C2B8-62-TN
%=== Câu 2 ===%
\begin{ex}
% ID: L12C2B8-62-TN
% Mức: TH
	{\color{red}\footnotesize\ttfamily [L12C1B1-TN-43]}\par %HienID
	Chuyển động Brown là bằng chứng thực nghiệm chứng tỏ điều gì về các phân tử chất lỏng, chất khí?
	\choice
	{Các phân tử đứng yên tại các vị trí cố định}
	{\True Các phân tử chuyển động hỗn loạn, không ngừng}
	{Các phân tử chỉ chuyển động khi có ánh sáng chiếu vào}
	{Các phân tử sắp xếp có trật tự chặt chẽ}
	\loigiai{
		Thí nghiệm của Brown quan sát thấy hạt phấn hoa lơ lửng trong nước chuyển động hỗn loạn không ngừng dù không chịu tác động bên ngoài nào khác, điều này chứng tỏ các phân tử nước xung quanh hạt phấn hoa luôn chuyển động hỗn loạn, không ngừng và va chạm vào hạt phấn hoa từ mọi phía.
		Kết luận: Các phân tử chuyển động hỗn loạn, không ngừng.
	}
\end{ex}

% ID: L12C1B1-TN-44
% Mức: NB
% ID: L12C1B1-93-TN
% ID: L12C1B1-93-TN
% Mức: NB
% ID: L12C2B8-63-TN
%=== Câu 3 ===%
\begin{ex}
% ID: L12C2B8-63-TN
% Mức: NB
	{\color{red}\footnotesize\ttfamily [L12C1B1-TN-44]}\par %HienID
	Khi nhiệt độ của một khối chất tăng lên, tốc độ chuyển động hỗn loạn trung bình của các phân tử cấu tạo nên chất đó sẽ
	\choice
	{giảm xuống}
	{không đổi}
	{\True tăng lên}
	{bằng không}
	\loigiai{
		Theo mô hình động học phân tử, nhiệt độ của vật càng cao thì tốc độ chuyển động (chuyển động nhiệt) của các phân tử cấu tạo nên vật càng lớn.
		Kết luận: Tốc độ chuyển động hỗn loạn trung bình của các phân tử tăng lên.
	}
\end{ex}

% ID: L12C1B1-TN-45
% Mức: NB
% ID: L12C1B1-94-TN
% ID: L12C1B1-94-TN
% Mức: NB
% ID: L12C2B8-64-TN
%=== Câu 4 ===%
\begin{ex}
% ID: L12C2B8-64-TN
% Mức: NB
	{\color{red}\footnotesize\ttfamily [L12C1B1-TN-45]}\par %HienID
	Trong thí nghiệm của Brown (năm 1827), đối tượng được quan sát trực tiếp dưới kính hiển vi là
	\choice
	{các phân tử nước}
	{\True các hạt phấn hoa rất nhỏ lơ lửng trong nước}
	{các ion hòa tan trong nước}
	{các phân tử khí hòa tan trong nước}
	\loigiai{
		Brown đã quan sát các hạt phấn hoa rất nhỏ lơ lửng trong nước bằng kính hiển vi và thấy chúng chuyển động hỗn loạn, không ngừng — hiện tượng này về sau được gọi là chuyển động Brown. Bản thân các phân tử nước có kích thước quá nhỏ để có thể quan sát trực tiếp bằng kính hiển vi quang học thời đó.
		Kết luận: Đối tượng quan sát trực tiếp là các hạt phấn hoa rất nhỏ lơ lửng trong nước.
	}
\end{ex}

% ID: L12C1B1-TN-2
% Mức: NB
% ID: L12C1B1-95-TN
% ID: L12C1B1-95-TN
% Mức: NB
% ID: L12C2B8-65-TN
%=== Câu 5 ===%
\begin{ex}
% ID: L12C2B8-65-TN
% Mức: NB
	{\color{red}\footnotesize\ttfamily [L12C1B1-TN-2]}\par %HienID
	Phát biểu nào \textbf{không} đúng với mô hình động học phân tử?
	\choice
	{Giữa các phân tử có lực tương tác gọi là lực liên kết phân tử}
	{\True Tốc độ chuyển động của các phân tử cấu tạo nên vật càng lớn thì thể tích của vật càng lớn}
	{Các chất được cấu tạo từ các hạt riêng biệt là phân tử}
	{Các phân tử chuyển động không ngừng}
	\loigiai{
		\begin{itemize}
			\item Tốc độ chuyển động của các phân tử cấu tạo nên vật càng lớn thì nhiệt độ càng lớn.
	\end{itemize}}
\end{ex}

% ID: L12C1B1-TN-17
% Mức: NB
% ID: L12C1B1-96-TN
% ID: L12C1B1-96-TN
% Mức: NB
% ID: L12C2B8-66-TN
%=== Câu 6 ===%
\begin{ex}
% ID: L12C2B8-66-TN
% Mức: NB
	{\color{red}\footnotesize\ttfamily [L12C1B1-TN-17]}\par %HienID
	Theo mô hình động học phân tử về cấu tạo chất thì các chất được cấu tạo từ các hạt riêng biệt là
	\choice
	{ion}
	{plasma}
	{nguyên tử}
	{\True phân tử}
	\loigiai{Vật chất được cấu tạo từ bởi 1 số rất lớn những hạt có kích thước rất nhỏ gọi là phân tử}
\end{ex}

% ID: L12C1B1-DS-13
% Mức: TH
% ID: L12C1B1-136-DS
% ID: L12C1B1-136-DS
% Mức: TH
% ID: L12C2B8-67-DS
%=== Câu 49 ===%
\begin{ex}
% ID: L12C2B8-67-DS
% Mức: TH
	{\color{red}\footnotesize\ttfamily [L12C1B1-DS-13]}\par %HienID
	Xét các phát biểu về mô hình động học phân tử và chuyển động Brown.
	\choiceTF
	{Mô hình động học phân tử cho rằng vật chất được cấu tạo liên tục, không có khoảng trống giữa các hạt}
	{\True Chuyển động Brown là bằng chứng thực nghiệm cho thấy các phân tử chất lỏng (hoặc khí) chuyển động hỗn loạn không ngừng}
	{\True Nhiệt độ của môi trường càng cao thì chuyển động Brown của hạt quan sát được càng mạnh (nhanh) hơn}
	{Chuyển động Brown chỉ xảy ra đối với các hạt có kích thước bằng đúng kích thước của một phân tử}
	\loigiai{
		\begin{itemchoice}
			\item Mô hình động học phân tử cho rằng vật chất được cấu tạo từ các hạt riêng biệt (gián đoạn), giữa các hạt luôn có khoảng cách chứ không phải cấu tạo liên tục. Mệnh đề sai.
			\item Hiện tượng hạt phấn hoa chuyển động hỗn loạn, không ngừng trong thí nghiệm của Brown chỉ có thể giải thích được nếu thừa nhận các phân tử chất lỏng xung quanh cũng đang chuyển động hỗn loạn, không ngừng và va chạm vào hạt phấn hoa. Mệnh đề đúng.
			\item Khi nhiệt độ môi trường tăng, tốc độ chuyển động nhiệt của các phân tử chất lỏng tăng lên, làm tần suất và cường độ va chạm vào hạt quan sát tăng, khiến chuyển động Brown của hạt diễn ra càng mạnh (nhanh) hơn. Mệnh đề đúng.
			\item Chuyển động Brown thường được quan sát ở các hạt có kích thước lớn hơn nhiều so với phân tử (như hạt phấn hoa), do sự mất cân bằng trong va chạm của rất nhiều phân tử môi trường lên hạt đó. Nếu hạt có kích thước quá lớn thì các va chạm sẽ triệt tiêu lẫn nhau và không quan sát được chuyển động hỗn loạn rõ rệt. Mệnh đề sai.
		\end{itemchoice}
	}
\end{ex}

% ID: L12C1B1-DS-14
% Mức: TH
% ID: L12C1B1-137-DS
% ID: L12C1B1-137-DS
% Mức: TH
% ID: L12C2B8-68-DS
%=== Câu 50 ===%
\begin{ex}
% ID: L12C2B8-68-DS
% Mức: TH
	{\color{red}\footnotesize\ttfamily [L12C1B1-DS-14]}\par %HienID
	Xét các phát biểu về lực liên kết phân tử.
	\choiceTF
	{\True Giữa các phân tử luôn tồn tại đồng thời cả lực hút và lực đẩy}
	{Khi khoảng cách giữa hai phân tử tăng lên, lực liên kết giữa chúng luôn tăng theo}
	{\True Lực liên kết phân tử là nguyên nhân giúp vật rắn giữ được hình dạng và thể tích xác định}
	{Ở thể khí, lực liên kết phân tử rất mạnh do khoảng cách giữa các phân tử khí rất lớn}
	\loigiai{
		\begin{itemchoice}
			\item Theo mô hình động học phân tử, giữa các phân tử luôn có đồng thời lực hút và lực đẩy; lực tổng hợp (lực liên kết) phụ thuộc vào khoảng cách giữa các phân tử. Mệnh đề đúng.
			\item Khi khoảng cách giữa hai phân tử tăng lên, lực liên kết giữa chúng giảm đi (yếu dần) chứ không tăng; ở khoảng cách đủ lớn, lực tương tác coi như không đáng kể. Mệnh đề sai.
			\item Nhờ lực liên kết phân tử đủ mạnh trong chất rắn (khoảng cách phân tử rất nhỏ), các phân tử chỉ dao động quanh vị trí cân bằng cố định, giúp vật rắn giữ được cả hình dạng và thể tích xác định. Mệnh đề đúng.
			\item Ở thể khí, khoảng cách trung bình giữa các phân tử rất lớn, mà lực liên kết phân tử giảm rất nhanh khi khoảng cách tăng, nên lực liên kết phân tử ở thể khí rất yếu chứ không mạnh. Mệnh đề sai.
		\end{itemchoice}
	}
\end{ex}

% ID: L12C1B1-TLN-1
% Mức: TH
% ID: L12C1B1-149-TLN
% ID: L12C1B1-149-TLN
% Mức: TH
% ID: L12C2B8-69-TLN
%=== Câu 62 ===%
\begin{bt}
% ID: L12C2B8-69-TLN
% Mức: TH
	{\color{red}\footnotesize\ttfamily [L12C1B1-TLN-1]}\par %HienID
	Cho các phát biểu sau về mô hình động học phân tử:\\
	(1) Vật chất được cấu tạo từ các hạt riêng biệt gọi là phân tử.\\
	(2) Các phân tử luôn đứng yên tại một vị trí cố định.\\
	(3) Nhiệt độ của vật càng cao thì tốc độ chuyển động của các phân tử càng lớn.\\
	(4) Giữa các phân tử không tồn tại lực tương tác nào.\\
	Có bao nhiêu phát biểu đúng?
	\shortans[oly]{$2$}
	\loigiai{
		Phát biểu (1) và (3) đúng theo nội dung mô hình động học phân tử. Phát biểu (2) sai vì các phân tử luôn chuyển động không ngừng. Phát biểu (4) sai vì giữa các phân tử luôn có lực liên kết (hút và đẩy). Vậy có $2$ phát biểu đúng.
	}
\end{bt}

% ID: L12C1B1-TLN-2
% Mức: TH
% ID: L12C1B1-150-TLN
% ID: L12C1B1-150-TLN
% Mức: TH
% ID: L12C2B8-70-TLN
%=== Câu 63 ===%
\begin{bt}
% ID: L12C2B8-70-TLN
% Mức: TH
	{\color{red}\footnotesize\ttfamily [L12C1B1-TLN-2]}\par %HienID
	Cho các phát biểu sau:\\
	(1) Chuyển động Brown là chuyển động của chính các phân tử chất lỏng.\\
	(2) Chuyển động Brown là chuyển động hỗn loạn của các hạt nhỏ lơ lửng trong chất lỏng do bị các phân tử chất lỏng va chạm.\\
	(3) Chuyển động Brown càng mạnh khi nhiệt độ môi trường càng cao.\\
	(4) Chuyển động Brown không phụ thuộc vào nhiệt độ.\\
	Có bao nhiêu phát biểu đúng?
	\shortans[oly]{$2$}
	\loigiai{
		Phát biểu (2) và (3) đúng. Phát biểu (1) sai vì chuyển động Brown là chuyển động của hạt lơ lửng (như hạt phấn hoa), không phải của chính phân tử chất lỏng. Phát biểu (4) sai vì chuyển động Brown càng mạnh khi nhiệt độ càng cao. Vậy có $2$ phát biểu đúng.
	}
\end{bt}

% ID: L12C1B1-TL-1
% Mức: VD
% ID: L12C1B1-160-TL
% ID: L12C1B1-160-TL
% Mức: VD
% ID: L12C2B8-71-TL
%=== Câu 73 ===%
\begin{bt}
% ID: L12C2B8-71-TL
% Mức: VD
	{\color{red}\footnotesize\ttfamily [L12C1B1-TL-1]}\par %HienID
	\begin{enumerate}
		\item Trình bày nội dung cơ bản của mô hình động học phân tử về cấu tạo chất.
		\item Giải thích tại sao khi nhỏ một giọt mực vào cốc nước, sau một thời gian mực sẽ loang ra khắp cốc nước dù không có sự khuấy trộn nào.
		\item Nếu tăng nhiệt độ của nước, hiện tượng loang mực ở ý 2) diễn ra nhanh hơn hay chậm hơn? Giải thích.
	\end{enumerate}
	\loigiai{
		\begin{enumerate}
			\item Mô hình động học phân tử về cấu tạo chất gồm ba nội dung cơ bản: (1) các chất được cấu tạo từ các hạt riêng biệt là phân tử; (2) các phân tử chuyển động không ngừng, nhiệt độ của vật càng cao thì tốc độ chuyển động của các phân tử càng lớn; (3) giữa các phân tử có lực hút và đẩy gọi chung là lực liên kết phân tử.
			\item Do các phân tử nước và phân tử mực đều chuyển động hỗn loạn không ngừng (chuyển động nhiệt), chúng va chạm và len lỏi xen lẫn vào nhau. Qua thời gian, các phân tử mực dần dần phân tán ra khắp thể tích nước, tạo nên hiện tượng loang mực (đây thực chất là hiện tượng khuếch tán, có cùng bản chất với chuyển động Brown).
			\item Hiện tượng loang mực sẽ diễn ra nhanh hơn. Khi nhiệt độ tăng, tốc độ chuyển động hỗn loạn của các phân tử nước và phân tử mực đều tăng lên, làm tần suất va chạm và tốc độ khuếch tán của các phân tử mực vào nước tăng lên, nên quá trình loang mực diễn ra nhanh hơn.
		\end{enumerate}
	}
\end{bt}

% ===== Câu 26 =====
% ID: L12C2B8-72-TL
% Mức: NB
\begin{bt}
% ID: L12C2B8-72-TL
% Mức: NB
Trình bày cách nhận biết và phương pháp giải dạng “Mô hình động học phân tử về cấu tạo chất và chuyển động Brown”. Phân tích nhận định sau và sửa lại nếu sai: “Các phân tử đứng yên hoàn toàn khi vật ở thể rắn.”
\loigiai{
Trước hết xác định hiện tượng, trạng thái và các đại lượng đã cho. Nêu định nghĩa hoặc hệ thức phù hợp, đổi về đơn vị SI, áp dụng đúng điều kiện và kiểm tra ý nghĩa kết quả. Nhận định cần đối chiếu: Các chất được cấu tạo từ các hạt riêng biệt; các hạt chuyển động không ngừng. Vì vậy nhận định đã cho sai; phát biểu đúng là: Các chất được cấu tạo từ các hạt riêng biệt; các hạt chuyển động không ngừng.
}
\end{bt}

% ===== Câu 27 =====
% ID: L12C2B8-73-TLN
% Mức: NB
\begin{ex}
% ID: L12C2B8-73-TLN
% Mức: NB
Cho bốn nhận định: (1) Các chất được cấu tạo từ các hạt riêng biệt; các hạt chuyển động không ngừng. (2) Các phân tử đứng yên hoàn toàn khi vật ở thể rắn. (3) Cần kiểm tra đơn vị và điều kiện áp dụng khi xử lí dạng “Mô hình động học phân tử về cấu tạo chất và chuyển động Brown”. (4) Có thể thay tùy ý nhiệt độ Celsius cho nhiệt độ tuyệt đối trong mọi công thức. Có bao nhiêu nhận định đúng? Trả lời bằng một số nguyên.
\shortans{2}
\loigiai{
Nhận định (1) và (3) đúng; (2) và (4) sai. Vậy có 2 nhận định đúng.
}
\end{ex}

% ===== Câu 29 =====
% ID: L12C2B8-74-DS
% Mức: NB
\begin{ex}
% ID: L12C2B8-74-DS
% Mức: NB
Xét các phát biểu sau về dạng “Mô hình động học phân tử về cấu tạo chất và chuyển động Brown” (bộ số 4).
\choiceTF
{\True Các chất được cấu tạo từ các hạt riêng biệt; các hạt chuyển động không ngừng.}
{Các phân tử đứng yên hoàn toàn khi vật ở thể rắn.}
{\True Khi giải dạng “Mô hình động học phân tử về cấu tạo chất và chuyển động Brown”, cần xác định đúng trạng thái đầu, trạng thái cuối và quy ước dấu hoặc điều kiện áp dụng.}
{Có thể bỏ qua mọi dữ kiện về khối lượng, nhiệt độ và hiệu suất trong mọi bài toán thuộc dạng này.}
\loigiai{
\begin{itemchoice}
\item (Đúng) Các chất được cấu tạo từ các hạt riêng biệt; các hạt chuyển động không ngừng. (Vì): Các chất được cấu tạo từ các hạt riêng biệt; các hạt chuyển động không ngừng. b) S: Các phân tử đứng yên hoàn toàn khi vật ở thể rắn. c) Đ: Khi giải dạng “Mô hình động học phân tử về cấu tạo chất và chuyển động Brown”, cần xác định đúng trạng thái đầu, trạng thái cuối và quy ước dấu hoặc điều kiện áp dụng. d) S: Có thể bỏ qua mọi dữ kiện về khối lượng, nhiệt độ và hiệu suất trong mọi bài toán thuộc dạng này. Kết luận dựa trên định nghĩa, điều kiện áp dụng và quy ước trong SGK KNTT
\item (Sai) Các phân tử đứng yên hoàn toàn khi vật ở thể rắn.
\item (Đúng) Khi giải dạng “Mô hình động học phân tử về cấu tạo chất và chuyển động Brown”, cần xác định đúng trạng thái đầu, trạng thái cuối và quy ước dấu hoặc điều kiện áp dụng.
\item (Sai) Có thể bỏ qua mọi dữ kiện về khối lượng, nhiệt độ và hiệu suất trong mọi bài toán thuộc dạng này.
\end{itemchoice}
}
\end{ex}

% ===== Câu 33 =====
% ID: L12C2B8-75-DS
% Mức: TH
\begin{ex}
% ID: L12C2B8-75-DS
% Mức: TH
Xét các phát biểu sau về dạng “Mô hình động học phân tử về cấu tạo chất và chuyển động Brown” (bộ số 1).
\choiceTF
{Các phân tử đứng yên hoàn toàn khi vật ở thể rắn.}
{\True Khi giải dạng “Mô hình động học phân tử về cấu tạo chất và chuyển động Brown”, cần xác định đúng trạng thái đầu, trạng thái cuối và quy ước dấu hoặc điều kiện áp dụng.}
{Có thể bỏ qua mọi dữ kiện về khối lượng, nhiệt độ và hiệu suất trong mọi bài toán thuộc dạng này.}
{\True Các chất được cấu tạo từ các hạt riêng biệt; các hạt chuyển động không ngừng.}
\loigiai{
\begin{itemchoice}
\item (Sai) Các phân tử đứng yên hoàn toàn khi vật ở thể rắn. (Vì): Các phân tử đứng yên hoàn toàn khi vật ở thể rắn. b) Đ: Khi giải dạng “Mô hình động học phân tử về cấu tạo chất và chuyển động Brown”, cần xác định đúng trạng thái đầu, trạng thái cuối và quy ước dấu hoặc điều kiện áp dụng. c) S: Có thể bỏ qua mọi dữ kiện về khối lượng, nhiệt độ và hiệu suất trong mọi bài toán thuộc dạng này. d) Đ: Các chất được cấu tạo từ các hạt riêng biệt; các hạt chuyển động không ngừng. Kết luận dựa trên định nghĩa, điều kiện áp dụng và quy ước trong SGK KNTT
\item (Đúng) Khi giải dạng “Mô hình động học phân tử về cấu tạo chất và chuyển động Brown”, cần xác định đúng trạng thái đầu, trạng thái cuối và quy ước dấu hoặc điều kiện áp dụng.
\item (Sai) Có thể bỏ qua mọi dữ kiện về khối lượng, nhiệt độ và hiệu suất trong mọi bài toán thuộc dạng này.
\item (Đúng) Các chất được cấu tạo từ các hạt riêng biệt; các hạt chuyển động không ngừng.
\end{itemchoice}
}
\end{ex}

% ===== Câu 34 =====
% ID: L12C2B8-76-TL
% Mức: TH
\begin{bt}
% ID: L12C2B8-76-TL
% Mức: TH
Trình bày cách nhận biết và phương pháp giải dạng “Mô hình động học phân tử về cấu tạo chất và chuyển động Brown”. Phân tích nhận định sau và sửa lại nếu sai: “Các chất được cấu tạo từ các hạt riêng biệt; các hạt chuyển động không ngừng.”
\loigiai{
Trước hết xác định hiện tượng, trạng thái và các đại lượng đã cho. Nêu định nghĩa hoặc hệ thức phù hợp, đổi về đơn vị SI, áp dụng đúng điều kiện và kiểm tra ý nghĩa kết quả. Nhận định cần đối chiếu: Các chất được cấu tạo từ các hạt riêng biệt; các hạt chuyển động không ngừng. Vì vậy nhận định đã cho đúng.
}
\end{bt}

% ===== Câu 37 =====
% ID: L12C2B8-77-DS
% Mức: TH
\begin{ex}
% ID: L12C2B8-77-DS
% Mức: TH
Xét các phát biểu sau về dạng “Mô hình động học phân tử về cấu tạo chất và chuyển động Brown” (bộ số 5).
\choiceTF
{Các phân tử đứng yên hoàn toàn khi vật ở thể rắn.}
{\True Khi giải dạng “Mô hình động học phân tử về cấu tạo chất và chuyển động Brown”, cần xác định đúng trạng thái đầu, trạng thái cuối và quy ước dấu hoặc điều kiện áp dụng.}
{Có thể bỏ qua mọi dữ kiện về khối lượng, nhiệt độ và hiệu suất trong mọi bài toán thuộc dạng này.}
{\True Các chất được cấu tạo từ các hạt riêng biệt; các hạt chuyển động không ngừng.}
\loigiai{
\begin{itemchoice}
\item (Sai) Các phân tử đứng yên hoàn toàn khi vật ở thể rắn. (Vì): Các phân tử đứng yên hoàn toàn khi vật ở thể rắn. b) Đ: Khi giải dạng “Mô hình động học phân tử về cấu tạo chất và chuyển động Brown”, cần xác định đúng trạng thái đầu, trạng thái cuối và quy ước dấu hoặc điều kiện áp dụng. c) S: Có thể bỏ qua mọi dữ kiện về khối lượng, nhiệt độ và hiệu suất trong mọi bài toán thuộc dạng này. d) Đ: Các chất được cấu tạo từ các hạt riêng biệt; các hạt chuyển động không ngừng. Kết luận dựa trên định nghĩa, điều kiện áp dụng và quy ước trong SGK KNTT
\item (Đúng) Khi giải dạng “Mô hình động học phân tử về cấu tạo chất và chuyển động Brown”, cần xác định đúng trạng thái đầu, trạng thái cuối và quy ước dấu hoặc điều kiện áp dụng.
\item (Sai) Có thể bỏ qua mọi dữ kiện về khối lượng, nhiệt độ và hiệu suất trong mọi bài toán thuộc dạng này.
\item (Đúng) Các chất được cấu tạo từ các hạt riêng biệt; các hạt chuyển động không ngừng.
\end{itemchoice}
}
\end{ex}

% ===== Câu 41 =====
% ID: L12C2B8-78-DS
% Mức: VD
\begin{ex}
% ID: L12C2B8-78-DS
% Mức: VD
Xét các phát biểu sau về dạng “Mô hình động học phân tử về cấu tạo chất và chuyển động Brown” (bộ số 2).
\choiceTF
{\True Khi giải dạng “Mô hình động học phân tử về cấu tạo chất và chuyển động Brown”, cần xác định đúng trạng thái đầu, trạng thái cuối và quy ước dấu hoặc điều kiện áp dụng.}
{Có thể bỏ qua mọi dữ kiện về khối lượng, nhiệt độ và hiệu suất trong mọi bài toán thuộc dạng này.}
{\True Các chất được cấu tạo từ các hạt riêng biệt; các hạt chuyển động không ngừng.}
{Các phân tử đứng yên hoàn toàn khi vật ở thể rắn.}
\loigiai{
\begin{itemchoice}
\item (Đúng) Khi giải dạng “Mô hình động học phân tử về cấu tạo chất và chuyển động Brown”, cần xác định đúng trạng thái đầu, trạng thái cuối và quy ước dấu hoặc điều kiện áp dụng. (Vì): Khi giải dạng “Mô hình động học phân tử về cấu tạo chất và chuyển động Brown”, cần xác định đúng trạng thái đầu, trạng thái cuối và quy ước dấu hoặc điều kiện áp dụng. b) S: Có thể bỏ qua mọi dữ kiện về khối lượng, nhiệt độ và hiệu suất trong mọi bài toán thuộc dạng này. c) Đ: Các chất được cấu tạo từ các hạt riêng biệt; các hạt chuyển động không ngừng. d) S: Các phân tử đứng yên hoàn toàn khi vật ở thể rắn. Kết luận dựa trên định nghĩa, điều kiện áp dụng và quy ước trong SGK KNTT
\item (Sai) Có thể bỏ qua mọi dữ kiện về khối lượng, nhiệt độ và hiệu suất trong mọi bài toán thuộc dạng này.
\item (Đúng) Các chất được cấu tạo từ các hạt riêng biệt; các hạt chuyển động không ngừng.
\item (Sai) Các phân tử đứng yên hoàn toàn khi vật ở thể rắn.
\end{itemchoice}
}
\end{ex}

% ===== Câu 45 =====
% ID: L12C2B8-79-DS
% Mức: VDC
\begin{ex}
% ID: L12C2B8-79-DS
% Mức: VDC
Xét các phát biểu sau về dạng “Mô hình động học phân tử về cấu tạo chất và chuyển động Brown” (bộ số 3).
\choiceTF
{Có thể bỏ qua mọi dữ kiện về khối lượng, nhiệt độ và hiệu suất trong mọi bài toán thuộc dạng này.}
{\True Các chất được cấu tạo từ các hạt riêng biệt; các hạt chuyển động không ngừng.}
{Các phân tử đứng yên hoàn toàn khi vật ở thể rắn.}
{\True Khi giải dạng “Mô hình động học phân tử về cấu tạo chất và chuyển động Brown”, cần xác định đúng trạng thái đầu, trạng thái cuối và quy ước dấu hoặc điều kiện áp dụng.}
\loigiai{
\begin{itemchoice}
\item (Sai) Có thể bỏ qua mọi dữ kiện về khối lượng, nhiệt độ và hiệu suất trong mọi bài toán thuộc dạng này. (Vì): Có thể bỏ qua mọi dữ kiện về khối lượng, nhiệt độ và hiệu suất trong mọi bài toán thuộc dạng này. b) Đ: Các chất được cấu tạo từ các hạt riêng biệt; các hạt chuyển động không ngừng. c) S: Các phân tử đứng yên hoàn toàn khi vật ở thể rắn. d) Đ: Khi giải dạng “Mô hình động học phân tử về cấu tạo chất và chuyển động Brown”, cần xác định đúng trạng thái đầu, trạng thái cuối và quy ước dấu hoặc điều kiện áp dụng. Kết luận dựa trên định nghĩa, điều kiện áp dụng và quy ước trong SGK KNTT
\item (Đúng) Các chất được cấu tạo từ các hạt riêng biệt; các hạt chuyển động không ngừng.
\item (Sai) Các phân tử đứng yên hoàn toàn khi vật ở thể rắn.
\item (Đúng) Khi giải dạng “Mô hình động học phân tử về cấu tạo chất và chuyển động Brown”, cần xác định đúng trạng thái đầu, trạng thái cuối và quy ước dấu hoặc điều kiện áp dụng.
\end{itemchoice}
}
\end{ex}

% ===== Câu 145 =====
% ID: L12C2B8-80-TN
% Mức: NB
\begin{ex}
% ID: L12C2B8-80-TN
% Mức: NB
Câu 1 thuộc dạng “Mô hình động học phân tử về cấu tạo chất và chuyển động Brown”. Phát biểu nào sau đây đúng?
\choice
{Các phân tử đứng yên hoàn toàn khi vật ở thể rắn.}
{Nhiệt lượng và nhiệt độ là hai đại lượng hoàn toàn giống nhau.}
{Mọi quá trình nhiệt học đều không phụ thuộc điều kiện thực hiện.}
{\True Các chất được cấu tạo từ các hạt riêng biệt; các hạt chuyển động không ngừng.}
\loigiai{
Phát biểu đúng là: Các chất được cấu tạo từ các hạt riêng biệt; các hạt chuyển động không ngừng. Các phương án còn lại trái với mô hình, định nghĩa hoặc hệ thức nhiệt học tương ứng.
}
\end{ex}

% ===== Câu 146 =====
% ID: L12C2B8-81-TN
% Mức: NB
\begin{ex}
% ID: L12C2B8-81-TN
% Mức: NB
Câu 5 thuộc dạng “Mô hình động học phân tử về cấu tạo chất và chuyển động Brown”. Phát biểu nào sau đây đúng?
\choice
{Các phân tử đứng yên hoàn toàn khi vật ở thể rắn.}
{Nhiệt lượng và nhiệt độ là hai đại lượng hoàn toàn giống nhau.}
{Mọi quá trình nhiệt học đều không phụ thuộc điều kiện thực hiện.}
{\True Các chất được cấu tạo từ các hạt riêng biệt; các hạt chuyển động không ngừng.}
\loigiai{
Phát biểu đúng là: Các chất được cấu tạo từ các hạt riêng biệt; các hạt chuyển động không ngừng. Các phương án còn lại trái với mô hình, định nghĩa hoặc hệ thức nhiệt học tương ứng.
}
\end{ex}

% ===== Câu 147 =====
% ID: L12C2B8-82-TN
% Mức: NB
\begin{ex}
% ID: L12C2B8-82-TN
% Mức: NB
Câu 9 thuộc dạng “Mô hình động học phân tử về cấu tạo chất và chuyển động Brown”. Phát biểu nào sau đây đúng?
\choice
{Các phân tử đứng yên hoàn toàn khi vật ở thể rắn.}
{Nhiệt lượng và nhiệt độ là hai đại lượng hoàn toàn giống nhau.}
{Mọi quá trình nhiệt học đều không phụ thuộc điều kiện thực hiện.}
{\True Các chất được cấu tạo từ các hạt riêng biệt; các hạt chuyển động không ngừng.}
\loigiai{
Phát biểu đúng là: Các chất được cấu tạo từ các hạt riêng biệt; các hạt chuyển động không ngừng. Các phương án còn lại trái với mô hình, định nghĩa hoặc hệ thức nhiệt học tương ứng.
}
\end{ex}

% ===== Câu 148 =====
% ID: L12C2B8-83-TN
% Mức: TH
\begin{ex}
% ID: L12C2B8-83-TN
% Mức: TH
Câu 2 thuộc dạng “Mô hình động học phân tử về cấu tạo chất và chuyển động Brown”. Phát biểu nào sau đây đúng?
\choice
{Nhiệt lượng và nhiệt độ là hai đại lượng hoàn toàn giống nhau.}
{Mọi quá trình nhiệt học đều không phụ thuộc điều kiện thực hiện.}
{\True Các chất được cấu tạo từ các hạt riêng biệt; các hạt chuyển động không ngừng.}
{Các phân tử đứng yên hoàn toàn khi vật ở thể rắn.}
\loigiai{
Phát biểu đúng là: Các chất được cấu tạo từ các hạt riêng biệt; các hạt chuyển động không ngừng. Các phương án còn lại trái với mô hình, định nghĩa hoặc hệ thức nhiệt học tương ứng.
}
\end{ex}

% ===== Câu 149 =====
% ID: L12C2B8-84-TN
% Mức: TH
\begin{ex}
% ID: L12C2B8-84-TN
% Mức: TH
Câu 6 thuộc dạng “Mô hình động học phân tử về cấu tạo chất và chuyển động Brown”. Phát biểu nào sau đây đúng?
\choice
{Nhiệt lượng và nhiệt độ là hai đại lượng hoàn toàn giống nhau.}
{Mọi quá trình nhiệt học đều không phụ thuộc điều kiện thực hiện.}
{\True Các chất được cấu tạo từ các hạt riêng biệt; các hạt chuyển động không ngừng.}
{Các phân tử đứng yên hoàn toàn khi vật ở thể rắn.}
\loigiai{
Phát biểu đúng là: Các chất được cấu tạo từ các hạt riêng biệt; các hạt chuyển động không ngừng. Các phương án còn lại trái với mô hình, định nghĩa hoặc hệ thức nhiệt học tương ứng.
}
\end{ex}

% ===== Câu 150 =====
% ID: L12C2B8-85-TN
% Mức: TH
\begin{ex}
% ID: L12C2B8-85-TN
% Mức: TH
Câu 10 thuộc dạng “Mô hình động học phân tử về cấu tạo chất và chuyển động Brown”. Phát biểu nào sau đây đúng?
\choice
{Nhiệt lượng và nhiệt độ là hai đại lượng hoàn toàn giống nhau.}
{Mọi quá trình nhiệt học đều không phụ thuộc điều kiện thực hiện.}
{\True Các chất được cấu tạo từ các hạt riêng biệt; các hạt chuyển động không ngừng.}
{Các phân tử đứng yên hoàn toàn khi vật ở thể rắn.}
\loigiai{
Phát biểu đúng là: Các chất được cấu tạo từ các hạt riêng biệt; các hạt chuyển động không ngừng. Các phương án còn lại trái với mô hình, định nghĩa hoặc hệ thức nhiệt học tương ứng.
}
\end{ex}

% ===== Câu 151 =====
% ID: L12C2B8-86-TN
% Mức: VD
\begin{ex}
% ID: L12C2B8-86-TN
% Mức: VD
Câu 3 thuộc dạng “Mô hình động học phân tử về cấu tạo chất và chuyển động Brown”. Phát biểu nào sau đây đúng?
\choice
{Mọi quá trình nhiệt học đều không phụ thuộc điều kiện thực hiện.}
{\True Các chất được cấu tạo từ các hạt riêng biệt; các hạt chuyển động không ngừng.}
{Các phân tử đứng yên hoàn toàn khi vật ở thể rắn.}
{Nhiệt lượng và nhiệt độ là hai đại lượng hoàn toàn giống nhau.}
\loigiai{
Phát biểu đúng là: Các chất được cấu tạo từ các hạt riêng biệt; các hạt chuyển động không ngừng. Các phương án còn lại trái với mô hình, định nghĩa hoặc hệ thức nhiệt học tương ứng.
}
\end{ex}

% ===== Câu 152 =====
% ID: L12C2B8-87-TN
% Mức: VD
\begin{ex}
% ID: L12C2B8-87-TN
% Mức: VD
Câu 7 thuộc dạng “Mô hình động học phân tử về cấu tạo chất và chuyển động Brown”. Phát biểu nào sau đây đúng?
\choice
{Mọi quá trình nhiệt học đều không phụ thuộc điều kiện thực hiện.}
{\True Các chất được cấu tạo từ các hạt riêng biệt; các hạt chuyển động không ngừng.}
{Các phân tử đứng yên hoàn toàn khi vật ở thể rắn.}
{Nhiệt lượng và nhiệt độ là hai đại lượng hoàn toàn giống nhau.}
\loigiai{
Phát biểu đúng là: Các chất được cấu tạo từ các hạt riêng biệt; các hạt chuyển động không ngừng. Các phương án còn lại trái với mô hình, định nghĩa hoặc hệ thức nhiệt học tương ứng.
}
\end{ex}

% ===== Câu 153 =====
% ID: L12C2B8-88-TN
% Mức: VDC
\begin{ex}
% ID: L12C2B8-88-TN
% Mức: VDC
Câu 4 thuộc dạng “Mô hình động học phân tử về cấu tạo chất và chuyển động Brown”. Phát biểu nào sau đây đúng?
\choice
{\True Các chất được cấu tạo từ các hạt riêng biệt; các hạt chuyển động không ngừng.}
{Các phân tử đứng yên hoàn toàn khi vật ở thể rắn.}
{Nhiệt lượng và nhiệt độ là hai đại lượng hoàn toàn giống nhau.}
{Mọi quá trình nhiệt học đều không phụ thuộc điều kiện thực hiện.}
\loigiai{
Phát biểu đúng là: Các chất được cấu tạo từ các hạt riêng biệt; các hạt chuyển động không ngừng. Các phương án còn lại trái với mô hình, định nghĩa hoặc hệ thức nhiệt học tương ứng.
}
\end{ex}

% ===== Câu 154 =====
% ID: L12C2B8-89-TN
% Mức: VDC
\begin{ex}
% ID: L12C2B8-89-TN
% Mức: VDC
Câu 8 thuộc dạng “Mô hình động học phân tử về cấu tạo chất và chuyển động Brown”. Phát biểu nào sau đây đúng?
\choice
{\True Các chất được cấu tạo từ các hạt riêng biệt; các hạt chuyển động không ngừng.}
{Các phân tử đứng yên hoàn toàn khi vật ở thể rắn.}
{Nhiệt lượng và nhiệt độ là hai đại lượng hoàn toàn giống nhau.}
{Mọi quá trình nhiệt học đều không phụ thuộc điều kiện thực hiện.}
\loigiai{
Phát biểu đúng là: Các chất được cấu tạo từ các hạt riêng biệt; các hạt chuyển động không ngừng. Các phương án còn lại trái với mô hình, định nghĩa hoặc hệ thức nhiệt học tương ứng.
}
\end{ex}
